\section{The remaining-exposure theorem}
\label{sec:exposure}

\begin{lemma}[Supersolution segment]\label{lem:segment}
Suppose $q\in(0,1]^m$, $\PhiB(q)\le0$, and let $h=1-q$. For $0\le a\le1$,
$q_a=1-ah$ satisfies $\PhiB(q_a)\le0$. In the binary case, exactly
\begin{equation}\label{eq:segment}
 \PhiB(1-ah)=a\PhiB(q)-a(1-a)\,b\odot D(h,h).
\end{equation}
If $Rq\le ch$ for $c\ge0$, then $Rq_a\le c(1-q_a)$.
\end{lemma}

\begin{proof}
Conservation of $Q$ and normalization $F(1)=1$ give $\PhiB(1)=0$. For binary
offspring, expand $D(1-ah,1-ah)$ and subtract $a\PhiB(1-h)$; the linear chemical
and death terms cancel as in \eqref{eq:segment}, and $D(h,h)\ge0$ because
$h\ge0$ and the daughter law has nonnegative coefficients. For general bounded
offspring, each monomial of $F_i(1-ah)$ is a product
$\prod_{\ell=1}^k(1-ah_{j_\ell})$, whose second derivative in the scalar $a$ is
a sum of terms each containing two nonnegative factors $h_{j}$ and otherwise
nonnegative factors; hence each monomial, and so $F_i(1-ah)$, is convex
\emph{as a function of the scalar} $a$ on $[0,1]$. Its graph lies below the
chord, so $F_i(1-ah)\le(1-a)F_i(1)+aF_i(q)$. The remaining terms are affine in
$a$, whence $\PhiB(1-ah)\le a\PhiB(q)\le0$. Finally $R(1-ah)=aRq$ because
$R\one=0$.
\end{proof}

\begin{remark}
The convexity used here is directional: it is convexity along the prescribed
segment $1-a(1-q)$, not joint convexity of the generating function on the cube.
Joint convexity is false in general --- the monomial $x_1x_2$ has an indefinite
Hessian --- and nonnegative mixed partial derivatives do not supply it. The
homotopy is affine in $a$; its later dependence on remaining exposure is
exponential.
\end{remark}

\begin{theorem}[Exposure floor under feedback]\label{thm:main}
Assume the model of Section~\ref{sec:model}. Let $q\in(0,1]^m$ and $c\ge0$
satisfy
\begin{equation}\label{eq:cert}
 \PhiB(q)\le0,\qquad Rq\le c(1-q).
\end{equation}
Suppose the declared continuation has extinction vector $q^{\mathrm{off}}\le q$.
Then for every admitted predictable policy and every initial configuration $z$,
\begin{equation}\label{eq:main}
 \Prob_z(|Z_T|>0)\ \ge\ r_\pi(T)\ \ge\
 1-\prod_i\left[1-e^{-cB}(1-q_i)\right]^{z_i}.
\end{equation}
The same conclusion holds at an almost surely finite adaptive withdrawal time
if the pathwise budget holds until that time and the same continuation
condition holds conditional on its history.
\end{theorem}

\begin{proof}
For remaining exposure $\beta\ge0$ set
\[
 a=e^{-c\beta},\qquad q_\beta=1-a(1-q),\qquad L(\beta,z)=1-V_{q_\beta}(z).
\]
All coordinates of $q_\beta$ are positive, even at $\beta=0$.
Lemma~\ref{lem:segment} gives $\Phi_v(q_\beta)\le cv(1-q_\beta)$ for every
admitted $v\ge0$, and $\partial_\beta q_\beta=c(1-q_\beta)$. Differentiating the
finite product in \eqref{eq:generator} yields
\begin{equation}\label{eq:barrier}
 (\mathcal G_v-v\partial_\beta)L(\beta,z)
 =V_{q_\beta}(z)\sum_i\frac{z_i}{q_{\beta,i}}
      \bigl[cv(1-q_{\beta,i})-\Phi_v(q_\beta)_i\bigr]\ \ge\ 0 .
\end{equation}
This also holds at $z=0$, where the sum is empty.

Let $\beta_t=B-C_t$, which is nonincreasing and nonnegative by
\eqref{eq:budget}. Stop at bounded population and proposal count and at time
$T$. The finite jump compensation identity and the finite-variation chain rule
make $L(\beta_t,Z_t)$ a submartingale after stopping. Non-explosion removes the
stopping sequence, and $0\le L\le1$ supplies dominated convergence; therefore
$\E L(\beta_T,Z_T)\ge L(B,z)$.

The function $L$ is nonincreasing in $\beta$. Conditional on the history at $T$,
the continuation survival probability is
$1-\prod_i(q_i^{\mathrm{off}})^{Z_i(T)}$, so pointwise
\[
 L(\beta_T,Z_T)\le L(0,Z_T)\le1-\prod_i(q_i^{\mathrm{off}})^{Z_i(T)}
 \le\one_{\{|Z_T|>0\}}.
\]
Taking expectations proves \eqref{eq:main}. Independence is used only for the
declared uncontrolled continuation, conditional on its starting population. For
an almost surely finite withdrawal time $\tau$, first stop at $\tau\wedge k$;
boundedness and eventual equality to the stopped process at $\tau$ allow
$k\to\infty$, after which the same terminal comparison applies. A rule that
waits for extinction is not admissible under this extension unless it is almost
surely finite.
\end{proof}

\begin{remark}
The hypothesis $\PhiB(q)\le0$ makes $q$ an upper bound for the \emph{minimal}
extinction fixed point: $V_q(Z_t)$ is a bounded supermartingale for the
uncontrolled source, extinction by time $t$ forces $V_q=1$, and letting
$t\to\infty$ gives $q^{\mathrm{off}}\le q$ whenever the continuation also
satisfies $\Phi_{\mathrm{off}}(q)\le0$. An arbitrary numerical root of
$\Phi=0$ cannot be substituted without identifying it as that minimal solution.
\end{remark}

\begin{corollary}[Preparations, actuators and uncertainty]\label{cor:extensions}
For a random initial configuration, average the right side of \eqref{eq:main}
over its actual law; for a preparation known only to lie in a set, minimize that
right side over the set. Suppose instead that
\[
 \Phi_u(x)=\PhiB(x)+\sum_ku_k\Psi_k(x),\qquad
 \Psi_k(x)=R_kx+\kappa_k\odot(1-x),\quad R_k\one=0,
\]
all admitted rates are valid, $u_k\ge0$, and $\Psi_k(q)\le c_k(1-q)$ with
$c_k\ge0$. Then replace $cB$ in \eqref{eq:main} by $\sum_kc_kB_k$ for pathwise
budgets $\int u_k\le B_k$. For a single cost budget $\int\sum_ka_ku_k\le B$ with
$c_k\le c_*a_k$, use $c_*B$. The statements are uniform over a fixed source
parameter set if the same certificates and continuation inequalities hold
throughout that set.
\end{corollary}

\begin{proof}
Condition on the preparation and apply Theorem~\ref{thm:main}. Linearity gives
$\Psi_k(q_a)=a\Psi_k(q)$, so the segment bound of Lemma~\ref{lem:segment} holds
with $\sum_kc_ku_k$ in place of $cv$. Use the remaining weighted exposure
$\beta=\sum_kc_kB_k-\int\sum_kc_ku_k$ and $a=e^{-\beta}$ in
\eqref{eq:barrier}. The cost-budget version uses
$\sum_kc_ku_k\le c_*\sum_ka_ku_k$. The proof holds for each fixed source
separately, giving uniformity without exchanging the policy and source
quantifiers.
\end{proof}

A zero $c_k$ is a property of this certificate; it does not prove that the
actuator is biologically useless. A mixture of initial populations must not be
replaced by its mean composition, since the right side of \eqref{eq:main} is not
affine in $z$. For application of a data-selected \emph{upper} risk certificate,
a joint confidence set with coverage $1-\alpha$ and a policy bounded by $\delta$
throughout that set gives unconditional deployment risk at most $\delta+\alpha$,
by conditioning on coverage and using independent deployment; independent
marginal intervals do not automatically provide the required joint coverage.

\begin{proposition}[Sharpness over the general class]\label{prop:sharp}
The functional form of \eqref{eq:main} is attained in a limiting admissible
comparison source. This does not by itself establish sharpness of the constants
in any particular molecular source.
\end{proposition}

\begin{proof}
Take one immortal baseline type, no reproduction, and controlled death intensity
$cv$. Every constant $q\in(0,1)$ has $\PhiB(q)=0$ and $\Psi(q)=c(1-q)$; after
withdrawal the source is immortal and $q^{\mathrm{off}}=0$. Under a
deterministic schedule using exposure $B$, independent survival of the $n$ cells
gives exactly $1-(1-e^{-cB})^n$. Letting $q\downarrow0$ in
Theorem~\ref{thm:main} recovers equality. No zero denominator is used.
\end{proof}

Proposition~\ref{prop:sharp} leaves open whether the form is attained inside a
source with genuine inheritance. Section~\ref{sec:combo} shows that it is, in
the following sense: in the molecular source with the added-death actuator the
ratio of certified sufficient to necessary exposure is bounded by $1.45$ as
$n/\delta\to\infty$, uniformly in the number of memory sites.
